\scanpage{290}
\begin{center}- 101 -\end{center}\hfill IX

envisagé dans (2.5.3), et plus précisément, nous verrons que sa connaissance équivaut essentiellement à celle de l'action du groupe de monodromie $I$ sur $U$.

9.2. \uline{Cas} $\ell\ne p$ . Plaçons nous donc d'abord dans le cas, plus simple, où $\ell\ne p$, en reprenant alors les notations introduites dans (2.5.3). Notons qu'à torsion près, la relation $V=U^I$ signifie aussi que $V^\perp$ est le module engendré par les éléments $gx-x$ ($u\in U$, $g\in I$), et qu'on a en tous cas
\[
(9.2.1)\qquad gx-x\in V^\perp\quad\text{pour }x\in U,\ g\in I\quad.
\]
On a donc, pour $g\in I$ :
\[
(9.2.2)\qquad gx=x+u(g).x\quad,
\]
où $g\longmapsto u(g)$ est un homomorphisme de $U$ dans $V^\perp$ qui s'annule sur $V=U^I$, ou ce qui revient au même
\[
(9.2.3)\qquad u\in\operatorname{Hom}(U/V,\ V^\perp)\simeq(U/V)^\vee\otimes V^\perp\quad.
\]
Dans le cas envisagé de la réduction semi-stable, i.e. quand on a $V^\perp\subset V$, la formule (9.2.2) exprime que l'opération $g$ dans $U$ est une \uline{transvection} parallèle à $V^\perp=W$, et la composition de ces transvections n'est autre que l'addition des homomorphismes $u$ correspondants (9.2.3). Donc l'opération du groupe d'inertie $I$ sur $U$ s'explicite entièrement alors par un homomorphisme de groupes
\[
(9.2.4)\qquad I\longrightarrow(U/V)^\vee\otimes V^\perp\quad.
\]
\margin{$\longleftarrow$}
\margin{\unsure{419}}
\margin{419}
\margin{291}
